\scanpage{25}
$M_{i_0}\in\mathrm{Ob}\,D^b(\mathcal M^+(X_{i_0}))$, d'où $M_i,M$, on % L1
aura pour $f_{i_0}:X_{i_0}-Y_{i_0}$ la relation % L2
\[
\mathbb R f_*(M)=\varinjlim_i v_i^*(\mathbb R f_{i*}(M_i)) % L3
\]
où $v_i:Y-Y_i$ est le morphisme % L4
canonique. % L5

4) Soit $\ell\in\mathbb P$ tel que $\ell\,1_X\in\Gamma(X,\mathscr O_X)$ soit % L6
inversible. Alors on a un foncteur % L7
\[
T_\ell=T_\ell^{(X)}:\mathcal M^+(X)\longrightarrow\mathcal M_\ell^+(X) % L8
\]
où $\mathcal M_\ell^{*}(X)$ est \inserted{la catégorie} \unsure{formé|formée} des ``$\mathbb Q_\ell$-modules % L9
constructibles sur $X$'', i.e. la catégorie % L10
déduite des la catégorie des ``systèmes % L11
$\ell$-adiques de faisceaux de $\ell$-torsion % L12
constructibles'' en négligeant précisément % L13
les faisceaux de torsion. Le foncteur % L14
$T_\ell$ est \inserted{\uuline{compatible avec} $\otimes$ \uuline{et unité}} \uuline{exact} et \fbox{\uuline{fidèle}} (mais non % L15
pleinement fidèle), \uuline{compatible avec le} % L16

% archive sheet number: 24
